% !TeX root = ../papers/nc-p3.tex
\section{固定楕円曲線上の二位置移動と四次障害の商}
\label{sec:fixed-exchange-isogeny}
前節の楕円曲線を $E_\sigma:S^2=F_\sigma(a)$ とし、
$s=S/(\lambda(z^2-1)(a^2-z))$ を $2px+m$ に対応させる。
\begin{theorem}\label{thm:fixed-exchange-isogeny}
通常倍率で二行列を交換する二位置移動は $E_\sigma$ 上の固定点のない対合である。
原点を選べば位数2の点による平行移動であり、商写像は次数2の同種写像となる。
次の行列を $M=\left(\begin{smallmatrix}A&B\\C&-A\end{smallmatrix}\right)$、
$d=\det M$, $D=Ca-A$, $b=(Aa+B)/D$ と書く。
\[
 M_+=\begin{pmatrix}2\lambda z&-z(\lambda^2z+1)\\
 \lambda^2z+1&-2\lambda z\end{pmatrix},
 \qquad M_-=\begin{pmatrix}0&-z\\1&0\end{pmatrix}.
\]
交換の式は $(a,S)\mapsto (b,dS/D^2)$ である。
$C\ne0$では商の関数体を次の四次モデルで表示できる。
\[
 U=a+b,\quad V=S\frac{dU}{da},\quad
 V^2=\left(U^2-4(AU+B)/C\right)H(U),\quad
 H(a+b)D^2=C^2F_\sigma(a).
\]
$H$は二次式である。$C=0$の場合にも商写像の次数と平行移動の記述は保たれる。
\end{theorem}
\begin{proof}
固定直線上では $g_1=\sigma z f_1$, $h_1=\sigma e_1$ で、負の直線では定数項が消える。
$2px+m=s$ と $2qy+h=-s$ を使うと、第二行列の非対角・対角成分は
\[
 f=(s-1)/(2\sigma z f_1),\qquad e=(1-s)/(2\sigma e_1).
\]
従って $e/f=-zf_1/e_1=b$ は符号に依存しない。交換して両行列を $f$ で割ると、
新しい第一行列の座標は $b$、第二の非対角成分は $1/f$ になる。
定理\ref{thm:kummer-polar-completion}のテンソル恒等式によりこれは二位置移動を実現する。

$M^2=-dI$で $d$ は正負の順に $z(\lambda^2z-1)^2,z$ となり、
$F_\sigma(b)D^4=d^2F_\sigma(a)$ である。符号も決定するため、
$f=f_0+\sigma f_1(a^2-zf^2)$ を用いる。これを法として
\[
 (b^2-z)D^2(f+2\sigma z f_1(b))
 =d(a^2-z) f(1+2\sigma z f_1(a)f)
\]
となるので、実際の行列交換は上記の正符号の持上げである。
\path{computations/verify_fixed_exchange_isogeny.py} は両符号で恒等式を検算する。
滑らかな射影曲線の双有理写像なので全域に延びる。$db/da=d/D^2$ で正則微分
$da/S$ を保つ。楕円曲線の自己同型は平行移動と原点を固定する自己同型の積であり、
標数零では後者の微分への作用は忠実である。基底上の変換が非自明な対合なので
非零の位数2の平行移動となる。

$C\ne0$で $a^2-Ua+(AU+B)/C=0$ となる。$F_\sigma/(a-A/C)^2$ は交換不変で、
$U$の無限遠にのみ高々二位の極を持つから二次式 $H(U)$ である。
$a$の無限遠での最高次係数は非零なので次数は正確に2である。
微分公式 $(a-A/C)dU/da=2a-U$ から $V^2$ の式が出る。
$V$は交換不変で、$S$は $V$ と $a$ から復元できる。従って生成した体は指数2である。
\end{proof}

\begin{corollary}
四始点のいずれかで長さ4の標準次元を失う条件と、標準Koszul左右AS正則である条件は、
交換が窓に作用する共通定義域でこの商に降りる。一般倍率では前節の十二の非正則表示は
二つの商曲線上の合計六点に写る。
\end{corollary}
\begin{proof}
交換は始点を二つ移し、各文字空間の基底を変える。従っていずれかの始点での次元増加も、
全始点に課す左右正則性も保つ。十二表示の表は
$(a,e,f)\mapsto(e/f,a/f,1/f)$で閉じる。一般倍率で表の点は相異なり、対合は自由である。
この商は障害判定のためであって、添字保持同型を与えるものとは限らない。
\end{proof}

具体的な階数検査も小さくできる。第 $i$ 辺の二次商を $Q_i$、関係を $R_i$ と書くと、
長さ4の道空間は
\[
 V_i\otimes R_{i+1}\otimes V_{i+3}\longrightarrow Q_i\otimes Q_{i+2}
\]
の余核である。外側二対で先に関係を割り、残る中央の挿入を写したからである。
四文字次数で標準窓の長さ4の次元は順に $(9,10,8,8)$ である。
実際、群環で $v$ を四元群の正則表現、$r=(2,2,1,1)$ を二次関係の文字とすれば、
Koszul列の三次・四次項はそれぞれ $v$ と自明文字 $1$ である。
長さ4の係数は $v^4-3v^2r+r^2+2v^2-1$ となる。上記の余域は
$(26,26,24,24)$ なので、標準値に対応する階数は $(17,16,16,16)$ である。

\path{computations/verify_fixed_quartic_ranks.py} は倍率 $(3,2)$、素体
$\mathbb F_{101}$ で $a$ が有限、$pqxy\det A\det B\ne0$ の有理点を全て調べる。
正の曲線の101点のうち97点、負の曲線の104点のうち96点で全始点が標準次元を持つ。
他は全て前節の表の還元で、増加する文字は常に自明文字である。
この有限体の列挙は、高次拡大上の点や標数零での障害因子の全体を決定しない。
